\section{专科基础模型的不可替代位置}

虽然通用基础模型不断刷新综合基准，但在结构化 EHR 时序、专科影像、组学预测等具备强结构先验的子领域，专科自监督基础模型在精度、样本效率与跨机构泛化三个维度上系统性优于通用多模态大模型。通用与专科并非替代关系，而更接近互补关系。

\subsection{结构化 EHR 时序：从 Med-BERT 到 Foresight 2}

结构化 EHR 是医疗人工智能中最具结构先验的子领域。它本质上是一组以患者为单位、按时间排序的 ICD、SNOMED、CPT 与药物编码序列，与自然语言相比具有三个关键结构特征：\textbf{编码空间离散且封闭、时间间隔关键、患者间方差远大于编码方差}。这些特征决定了直接套用通用 LLM 在 EHR 上很难取得最优表现。

Med-BERT 与 CEHR-BERT 分别代表结构先验工程化的两条早期路径[30,31]。Med-BERT 在 Cerner Health Facts 的 2,849 万患者诊断码序列上预训练 17M 参数的 BERT 风格模型，关键创新是在 masked LM 之外引入延长住院时长预测作为领域特定的第二预训练目标；在 300 个样本的小数据场景下，Med-BERT 接 Bi-GRU 仍可达 AUC 0.75，相当于训练集扩大约 10 倍后的 Bi-GRU 性能。CEHR-BERT 则聚焦 \textbf{时间结构}，以人工时间标记把相邻就诊间隔以离散 token 形式插入序列，配合概念、时间与年龄嵌入，并以就诊类型预测替换 BERT 原版 NSP；在仅用 5\% 训练数据微调的情况下，CEHR-BERT 仍能超越使用全量数据训练的对比模型。这两项工作共同揭示出：\textbf{样本效率会随结构先验显著放大}。

Foresight 2 则进一步把结构先验与通用语言能力结合起来[41]。其关键设计是 \textbf{语境化患者时间线}：保留每个 SNOMED 概念被提及的原始文本上下文，并把 SNOMED 编码嵌入 LLM 分词器，让 7B 规模的通用 LLM（Mistral-7B 与 LLaMAv2-7B）能够高效处理标准化医学本体的输入与输出。在 MIMIC-III 535 名患者的风险预测测试集上，Foresight 2 的 P@5 达到 0.90，显著高于 GPT-4-turbo 的 0.65。这一对比直接说明：当 EHR 任务被恰当地结构化建模时，\textbf{7B 专科模型可以系统性优于参数规模更大的通用模型}。同样重要的是，Foresight 2 的消融实验显示，移除患者时间线中的文本上下文后性能下降约 40\%，说明结构先验与通用语言能力必须同时具备。

\subsection{影像与分割专科基础模型}

专科影像领域的方法学创新沿三条线索展开：\textbf{视觉自监督表示学习、视觉语言对齐、参数高效适配}。这三条线索分别对应不同的下游任务结构与数据可得性约束，但都在小样本与长尾场景下展现出对通用模型的系统性优势。

视觉自监督一脉以 RETFound 为旗舰[32]。其以 ViT-Large 在 90.4 万张彩色眼底照片与 73.6 万张 OCT 图像上做 MAE 风格预训练，下游覆盖糖网分级、青光眼、AMD 等 8 项眼科诊断任务以及心衰、心梗、缺血性卒中、帕金森病等 4 项三年预后预测；全部 12 项任务的 AUROC 均优于 ImageNet 预训练与监督基线，且 \textbf{优势在每类样本少于 100 时最为显著}。Pai 等进一步把同一思路扩展到 3D CT，在跨机构外部验证集上 AUC 与 C-index 显著优于 ImageNet-3D 与从头训练基线[38]。

视觉语言对齐一脉的代表是 CONCH[37]。其方法学包含两步：先以 1,600 万张病理 ROI 做 iBOT 风格视觉自监督训练，再在 117 万对病理图像--文本上以 CoCa 框架做图文对比与描述生成联合训练。这种 \textbf{视觉表示与语言接地分层使用} 的设计让 CONCH 在 14 类基准上有 11 项达到 SOTA，且在 zero-shot 罕见癌种分类上较 PLIP 提升超过 10 个百分点。

参数高效适配一脉则以 SAM-Med2D 与 MedSAM 为代表[36,39]。二者均以 Meta SAM 为骨架，通过 adapter、continual pretraining 等手段做专科适配。SAM-Med2D 合并 31 个公开数据集，在 9 类外部数据集上 Dice 较原生 SAM 提升 10--30 个百分点；MedSAM 则在不同数据规模与训练协议下复现了相同思路。这条路线以通用基础模型为骨架、以参数高效手段做专科适配，在精度与经济性之间取得了现实平衡。

\subsection{通用模型局限性的对照证据}

Jin 等发表在 \textit{npj Digital Medicine} 的研究为通用多模态模型在同类任务上的局限提供了直接对照证据[45]。他们以多类医学影像任务评估 GPT-4V 的多模态推理能力，发现一个反复出现的失效模式：\textbf{模型的选项准确率并不低，但图像感知与推理路径中包含大量隐性错误}。也就是说，模型往往凭借选项线索给出正确答案，但其对图像的实际描述与图中内容不符，或者关键区域识别错误。由此可见，单看选择题准确率会高估通用多模态模型在临床部署中的可靠性。

\subsection{章节小结与发展图景}

Foresight 2 在 EHR 上压制 GPT-4-turbo，RETFound 在小样本眼科任务上优于 ImageNet-ViT，CONCH 在 zero-shot 罕见癌种任务上优于 PLIP，Pai 等在小样本预后任务上优于 ImageNet-3D。综上，在结构化、长尾、小样本三类特征同时出现的子领域，\textbf{专科自监督基础模型仍然不可替代}。

综合前文可见，能力扩张的来源是数据规模、训练目标与结构先验三者的乘积，没有任何一项能够单独决定上限。当下“通用多模态大模型 + 专科基础模型 + 检索与工具调用”的混合架构，是这一时期最具代表性的事实图景。

\begin{table}[htbp]
\centering
\caption{通用模型与专科基础模型在四类任务上的相对优势矩阵}
\label{tab:specialist-matrix}
\scriptsize
\begin{tabularx}{\textwidth}{L{2.3cm} L{2.5cm} Y Y L{2cm} L{1.6cm}}
\toprule
任务类别 & 代表性任务 & 通用大模型 & 专科基础模型 & 相对优势归属 & 证据来源 \\
\midrule
结构化时序预测 & EHR 风险预测、再入院、HF 预后 & GPT-4-turbo P@5 0.65 & Foresight 2 P@5 0.90；Med-BERT 小样本提升约 20\% AUC & 专科压制 & [41,30,31] \\
罕见类别长尾分类 & 罕见癌种 zero-shot、罕见病眼科 & PLIP、CLIP zero-shot 较弱 & CONCH zero-shot 提升超过 10 个百分点；RETFound 小样本 SOTA & 专科压制 & [37,32] \\
高分辨率影像几何 & 3D CT 报告、病理 WSI、跨模态分割 & GPT-4V 存在隐性感知失效；Med-Gemini-3D 幻觉明显 & CONCH、Pai 等、MedSAM 跨机构更稳健 & 专科压制 + 通用辅助 & [38,39,45] \\
开放语义多轮交互 & 多轮诊断对话、文献综述、患者教育 & AMIE / Med-PaLM 2 在 OSCE 等任务中占优 & 专科 LLM 缺失或较弱 & 通用压制 & [17,13] \\
\bottomrule
\end{tabularx}
\end{table}
